\chapter{Fundamentação teórica}

\section{Baterias de lítio em configuração 4S}

Uma bateria em configuração \textit{4S} possui quatro células em série. Nesse arranjo, a tensão
total do pack resulta da soma das tensões individuais das células, o que exige monitoramento
por célula para evitar condições de sobretensão e subtensão.

\section{Battery Management System}

Um BMS é responsável por monitorar variáveis elétricas e térmicas do pack, proteger a bateria,
gerenciar carga e descarga e manter a operação dentro de limites seguros \cite{andrea2010}.

\section{Proteção em BMS}

As proteções mais relevantes em BMS incluem:

\begin{itemize}
  \item sobretensão por célula;
  \item subtensão por célula;
  \item sobrecorrente de carga;
  \item sobrecorrente de descarga;
  \item curto-circuito;
  \item sobretemperatura e subtemperatura.
\end{itemize}

Segundo Andrea \cite{andrea2010}, a robustez da proteção depende de uma combinação de hardware
dedicado e lógica de controle, especialmente em sistemas multicélula.

\section{Balanceamento de células}

O balanceamento busca reduzir a divergência entre tensões individuais das células. Entre as
estratégias possíveis, o balanceamento passivo é uma das mais comuns em aplicações de menor
complexidade, por dissipar energia das células mais carregadas em resistores de bleed
\cite{plett2015v2}.

\section{Analog Front End em BMS}

O \textit{Analog Front End} é o bloco responsável pela leitura segura de tensões, correntes e
temperaturas, além de frequentemente incorporar mecanismos de detecção de falhas e suporte a
balanceamento. Em sistemas reais, o AFE reduz a dependência do ADC do microcontrolador e
melhora a confiabilidade da proteção \cite{ti_bq76920,ad_passive_balancing}.

\section{Microcontroladores na supervisão da BMS}

O microcontrolador pode assumir as tarefas de:

\begin{itemize}
  \item supervisão de alto nível;
  \item telemetria;
  \item máquina de estados;
  \item registro de falhas;
  \item comunicação com interfaces externas.
\end{itemize}

No contexto deste trabalho, o ESP32-S3 foi escolhido por sua capacidade de processamento,
interfaces de comunicação e potencial de conectividade sem fio.
